\documentclass[11pt]{article}

\usepackage[margin=1in]{geometry}
\usepackage{fontspec}
\IfFontExistsTF{Avenir Next}
  {\setmainfont{Avenir Next}}
  {\setmainfont{TeX Gyre Heros}}
\usepackage{amsmath,amssymb,mathtools}
\usepackage{booktabs}
\usepackage{enumitem}
\usepackage{microtype}
\usepackage{setspace}
\usepackage{xcolor}

\definecolor{navy}{HTML}{0C2852}
\definecolor{red}{HTML}{982A34}
\definecolor{bluewash}{HTML}{F5F6F8}
\definecolor{redwash}{HTML}{F9F5F5}

\setlength{\parindent}{0pt}
\setlength{\parskip}{0.55em}
\setstretch{1.5}
\setlist[itemize]{leftmargin=1.5em,itemsep=0.35em,topsep=0.3em}
\newcommand{\E}{\mathbb{E}}
\newcommand{\Var}{\operatorname{Var}}

\newenvironment{auditbox}
  {\begin{center}\begin{minipage}{0.94\textwidth}\color{navy}\hrule\vspace{0.55em}}
  {\vspace{0.45em}\hrule\end{minipage}\end{center}}

\begin{document}

\begin{titlepage}
  \vspace*{2.4cm}
  {\color{navy}\Huge\bfseries Variance Audit and Extensions\par}
  \vspace{0.45cm}
  {\Large\color{red} The Economics of Bicycles for the Mind\par}
  \vspace{0.65cm}
  {\color{red}\rule{\textwidth}{0.7pt}\par}
  \vspace{1.0cm}
  {\large Gabriel Saco\par}
  \vspace{0.15cm}
  {Universidad del Pac\'ifico\par}
  \vfill
  {\small Technical companion to Agrawal, Gans, and Goldfarb (2025)\par}
  \vspace{0.15cm}
  {\small August 24, 2026\par}
\end{titlepage}

\section*{Purpose and main finding}

This note checks Propositions 2 and 3 of Agrawal, Gans, and Goldfarb (2025).
The paper interprets continuation value, $V(\theta)$, as wages. Its absolute
inequality object is therefore the cross-sectional variance of total
continuation value, not the variance of an individual's adoption benefit.

\begin{auditbox}
\textbf{Audit conclusion.} Within the common interior region, total-value
variance can be U-shaped in continuous tool quality. The initial decline needs
the paper's heterogeneity condition, while a positive slope specifically at
$\theta=1$ needs the additional restriction $\theta^\star<1$. By contrast,
individual-benefit variance rises monotonically whenever $\Gamma/s$ is
heterogeneous.
\end{auditbox}

\section{Specialized solution and its domain}

Under the paper's specialized technology,
\begin{equation}
  p(se;\theta)=\sqrt{se+\theta},
  \qquad c(e)=e,
\end{equation}
the per-opportunity objective and first-order condition are
\begin{align}
  M(e;\theta)&=\alpha\Delta\sqrt{se+\theta}-e,\\
  \frac{\alpha\Delta s}{2\sqrt{se+\theta}}&=1.
\end{align}
Solving the first-order condition gives
\begin{align}
  e^*(\theta)&=\frac{\alpha^2\Delta^2s}{4}-\frac{\theta}{s},\\
  M(\theta)&=\frac{\alpha^2\Delta^2s}{4}+\frac{\theta}{s}.
\end{align}
The interior solution is valid only when
\begin{equation}
  \theta<\frac{\alpha^2\Delta^2s^2}{4}.
  \label{eq:interior}
\end{equation}

Let the discounted opportunity multiplier be
\begin{equation}
  \Gamma=\frac{\gamma_0}{1-\delta\gamma},
  \qquad K\equiv\frac{\Delta^2}{4}.
\end{equation}
Then continuation value inside the common interior region is
\begin{equation}
  V(\theta)=\Gamma\left(K\alpha^2s+\frac{\theta}{s}\right).
  \label{eq:value}
\end{equation}

The constrained optimum, which is also meaningful outside that region, is
\begin{equation}
  e^*(\theta)=\max\left\{K\alpha^2s-\frac{\theta}{s},0\right\}.
\end{equation}
At the corner, the per-opportunity value is
$M(0;\theta)=\alpha\Delta\sqrt{\theta}$. The quadratic variance formula below
must not be extrapolated through a population in which some workers have
reached this corner.

\section{Cross-sectional variance of continuation value}

Independence of $\Gamma$ from $(\alpha,s)$ implies
\begin{equation}
  \Var(V)
  =\E[\Gamma^2]\Var(M)+\Var(\Gamma)\E[M]^2.
  \label{eq:productvar}
\end{equation}
The first two moments of the per-opportunity value can be organized as
\begin{align}
  \E[M]&=K\E[\alpha^2]\mu_s+\theta\E[1/s],\\
  \Var(M)&=A+B\theta+C\theta^2,
\end{align}
where
\begin{align}
  A&=K^2\Var(\alpha^2s),\\
  B&=2K\E[\alpha^2]\left\{1-\mu_s\E[1/s]\right\},\\
  C&=\Var(1/s).
\end{align}
Differentiating equation~\eqref{eq:productvar} and collecting terms yields
\begin{equation}
  \frac{d\Var[V(\theta)]}{d\theta}
  =a_0+2\theta\Var(\Gamma/s),
  \label{eq:slope}
\end{equation}
with
\begin{equation}
  a_0=\frac{\Delta^2\E[\alpha^2]}{2}
  \left\{\E[\Gamma^2]
  -\E[\Gamma]^2\mu_s\E[1/s]\right\}.
  \label{eq:a0}
\end{equation}

Condition (30) in the paper is
\begin{equation}
  \frac{\E[\Gamma^2]}{\E[\Gamma]^2}
  <\mu_s\E[1/s].
  \label{eq:condition30}
\end{equation}
It is exactly the condition $a_0<0$. If, in addition,
$\Var(\Gamma/s)>0$, variance is strictly convex and has one minimum:
\begin{equation}
  \theta^\star
  =-\frac{a_0}{2\Var(\Gamma/s)}>0.
  \label{eq:turning}
\end{equation}
Thus condition (30) proves an initial decline followed by an eventual upturn
as long as the relevant range remains interior. It does not locate that
upturn before any particular positive value of tool quality.

\section{Individual adoption benefit is a different variance}

For worker $i$, the adoption gain relative to no tool is
\begin{equation}
  D_i(\theta)
  =V_i(\theta)-V_i(0)
  =\theta\frac{\Gamma_i}{s_i}.
\end{equation}
Therefore,
\begin{align}
  \Var[D(\theta)]&=\theta^2\Var(\Gamma/s),\\
  \frac{d\Var[D(\theta)]}{d\theta}
    &=2\theta\Var(\Gamma/s)>0
    \qquad (\theta>0).
\end{align}
This variance is monotone, not U-shaped. The economic distinction is simple:
total continuation value contains baseline heterogeneity as well as the tool
gain; the individual benefit isolates only the tool gain.

\section{Internal slips in the variance section}

\subsection{Equation (32) needs another condition}

Equation (32) states that condition (30) makes the variance slope negative at
$\theta=0$ and positive at $\theta=1$. The first statement follows from
equation~\eqref{eq:a0}. The second is equivalent to
\begin{equation}
  a_0+2\Var(\Gamma/s)>0
  \quad\Longleftrightarrow\quad
  \theta^\star<1,
\end{equation}
which is an additional restriction absent from the proposition.

\subsection{The printed positivity condition is insufficient}

The text before Proposition 3 assumes
$\Delta>2/(s\sqrt{\alpha})$ ``so that optimal effort is positive.'' At
$\theta=1$, however, the solved effort formula requires
\begin{equation}
  \Delta>\frac{2}{\alpha s}.
\end{equation}
For $0<\alpha<1$, the printed condition is weaker and does not guarantee
positive effort.

\subsection{Probability and interiority clash at one}

At an interior optimum,
\begin{equation}
  p(se^*;\theta)
  =\sqrt{se^*+\theta}
  =\frac{\alpha\Delta s}{2}.
\end{equation}
Positive effort at $\theta=1$ requires $\alpha\Delta s/2>1$, whereas the
general model defines $p\in[0,1]$. The square-root expression therefore cannot
be both a literal probability and an interior specification at that endpoint.
It must instead be read as an uncapped productivity index, capped or rescaled,
or solved with the corner.

\section{Two Proposition 2 caveats}

\subsection{General opportunity timing}

For
\begin{equation}
  \Gamma=\sum_{k=0}^{\infty}
  \left(\prod_{i=0}^{k}\gamma(i)\right)\delta^k,
\end{equation}
the derivative of the entire continuation multiplier is
\begin{equation}
  \frac{\partial\Gamma}{\partial\gamma(t)}
  =\sum_{k=t}^{\infty}\delta^k
  \prod_{\substack{i=0\\ i\ne t}}^{k}\gamma(i).
\end{equation}
The displayed ratio in Proposition 2 retains only the $k=t$ term. A general
early-versus-late ranking depends on the full opportunity path. The geometric
special case can still deliver the paper's intended discounting result.

\subsection{Tool-skill substitutability at the optimizer}

The envelope theorem gives
\begin{equation}
  \frac{\partial[V_0(1)-V_0(0)]}{\partial s}
  =\Gamma\alpha\Delta
  \left[p_s(se^*(1);1)-p_s(se^*(0);0)\right].
\end{equation}
The condition $p_{s\theta}<0$ compares tool qualities while holding effort
fixed. Because $e^*(1)\ne e^*(0)$, it does not by itself imply the
optimizer-level inequality inside the brackets. An additional restriction on
the effort effect is required.

\section{Exact-moment counterexample}

Let all variables be mutually independent and distributed as
\begin{equation}
  \alpha\sim U(.75,.85),\quad
  \gamma_0\sim U(.48,.52),\quad
  \gamma\sim U(.36,.44),\quad
  s\sim U(.8,1),
\end{equation}
with $\delta=.8$ and $\Delta=5$. Every variable has positive support and
satisfies the paper's restriction that its mean exceed three standard
deviations. Exact moments give
\begin{equation}
  1.0012732
  =\frac{\E[\Gamma^2]}{\E[\Gamma]^2}
  <\mu_s\E[1/s]
  =1.0041460.
\end{equation}
Thus condition (30) holds. Nevertheless,
\begin{align}
  \theta^\star&=1.7007141,\\
  \left.\frac{d\Var(V)}{d\theta}\right|_{\theta=1}
  &=-0.0051337<0.
\end{align}
The smallest interior boundary in the population is $2.25$, so the turning
point remains inside the common interior region. This is a counterexample to
the slope-at-one claim, not to the conditional U-shape itself.

\section{Limiting cases and corner behavior}

\begin{itemize}
  \item If $\Gamma/s$ is homogeneous, there is no strict quadratic curvature.
  \item With constant $\Gamma$ but heterogeneous skill, inverse skill bias can
    reduce total-value variance initially, while benefit variance still rises.
  \item With constant skill and heterogeneous $\Gamma$, the inverse-skill
    equalizing channel disappears and tool quality raises variance.
  \item At large $\theta$, heterogeneous workers reach zero effort at different
    thresholds. The interior U-shape is therefore not a global variance result.
\end{itemize}

\clearpage
\section*{Verdict}

The phrase ``variance is U-shaped in the level of AI'' is incomplete. The
paper studies the cross-sectional variance of total continuation value with
respect to continuous tool quality, subject to a heterogeneity threshold and a
common interior solution. Individual-benefit variance rises monotonically, and
the paper's endpoint claim requires the missing condition
$\theta^\star<1$.

\section*{Reference}

Agrawal, Ajay K., Joshua S. Gans, and Avi Goldfarb. 2025. ``The Economics of
Bicycles for the Mind.'' NBER Working Paper 34034.
doi:10.3386/w34034.

\end{document}
